% Compile from this directory with: pdflatex project_1.tex (twice).
\documentclass[11pt,a4paper]{article}
\usepackage[T1]{fontenc}
\usepackage[utf8]{inputenc}
\usepackage[margin=22mm]{geometry}
\usepackage{lmodern}
\usepackage{amsmath,amssymb}
\usepackage{enumitem}
\usepackage{xcolor}
\usepackage{microtype}
\usepackage{hyperref}
\definecolor{epflred}{HTML}{E3000B}
\hypersetup{colorlinks=true,urlcolor=epflred,linkcolor=epflred,
  pdftitle={ENG-654 Project 1: Identifying and Tracking the Eight IK Branches of the PUMA 560}}
\setlength{\parindent}{0pt}
\setlength{\parskip}{6pt}
\setlist[itemize]{leftmargin=*,itemsep=5pt,topsep=4pt}
\urlstyle{tt}

\begin{document}
{\small\textbf{EPFL \textbar\ ENG-654}\hfill Kinematics-Grounded Motion Planning for Robots}
\vspace{5mm}

{\color{epflred}\Large\bfseries Project 1}

{\LARGE\bfseries Identifying and Tracking the Eight IK Branches of the PUMA 560\par}

\textbf{Format:} Group of two students; simulation-based project.\\
\textbf{Course foundations:} Lectures 2--7 and Exercises 1--3.

\section*{Motivation and historical context}
A Cartesian path specifies where a robot's tool should move, but does not uniquely
specify the robot's posture. Although the PUMA 560 has six revolute joints and is
nonredundant for a full-pose task, a regular target can admit up to eight isolated
inverse-kinematic (IK) solutions. These postures reproduce the same tool pose but
can have different joint-limit margins, singularity proximity, and joint travel
along a subsequent path. Understanding their identity and continuation is therefore
essential to choosing an initial posture that can execute the whole motion.
Pointwise reachability alone does not establish path feasibility.

The PUMA family grew out of Victor Scheinman's electrically driven robot-arm work
and was developed at Unimation for industrial assembly in the late 1970s. The
Smithsonian dates its Unimate PUMA 500 to 1978; the Computer History Museum
describes the family's use in automobile assembly and other industrial
tasks.\footnote{\href{https://www.si.edu/object/unimate-puma-500:nmah_1028521}{Smithsonian: Unimate PUMA 500};
\href{https://www.computerhistory.org/timeline/ai-robotics/}{Computer History Museum: AI \& Robotics timeline}.}
The PUMA 560 became a widely studied example of a \emph{wrist-partitioned 6R}
manipulator: three arm joints position a wrist center, while three concurrent
wrist axes determine orientation. This spherical-wrist geometry permits the
full-pose IK problem to be separated into arm and wrist subproblems, yielding
the familiar shoulder, elbow, and wrist choices.\footnote{\href{https://www.petercorke.com/RTB/r9/html/SerialLink.html}{Peter Corke, Robotics Toolbox: SerialLink, analytical spherical-wrist IK}.}

Six-axis arms remain relevant to industrial assembly and handling; current
assembly applications include those documented by
St\"aubli.\footnote{\href{https://www.staubli.com/global/en/robotics/applications/assembly-robot.html}{St\"aubli: robotic assembly applications}.}
The planning questions studied here also arise when a tool must follow a process
contour with a prescribed orientation. The PUMA provides a tractable model for
studying these questions; its eight-class structure must not be assumed to apply
to every industrial 6R geometry.

\section*{Task}
\textbf{Overall objectives.} Build a consistent kinematic description of the PUMA
560, identify its eight regular IK classes, and evaluate whether this classification
helps disambiguate numerical solutions near singularities. Then determine how the
initial IK class affects the ability to follow a prescribed Cartesian pose path.

\begin{itemize}
  \item \textbf{Understand and validate the model.}
  Inspect the supplied URDF: link/joint tree, joint origins and axes, limits,
  fixed transforms, and visual mesh origins and scales. Check link-name and mesh-path
  consistency when importing it. Establish a D--H table and document joint-angle
  offsets and base/tool conventions. Compare D--H forward kinematics with the
  URDF-based model at several configurations using position and orientation
  residuals. Explain why moving an STL visual origin does not change kinematics.

  \item \textbf{Recover and study eight IK solutions.}
  Use a regular full-pose target admitting eight mathematical solutions, starting
  with the Lecture 4 example. Reuse or adapt the course analytical method or your
  validated Exercise 2 implementation. Visualize every solution and verify pose
  closure. Distinguish mathematical solutions from those admitted by the chosen
  joint limits; record angle periodicity and remove equivalent duplicates.

  \item \textbf{Develop a geometric classifier.}
  Assign shoulder (left/right), elbow (up/down), and wrist (flip/non-flip) labels.
  Derive identification rules consistent with your D--H conventions and relate
  them to the singularity factors in Lecture 6:
  \[
    \det J(q)=K f_S(q)f_E(q)f_W(q),\qquad
    (S,E,W)=\bigl(\operatorname{sgn}f_S,
    \operatorname{sgn}f_E,\operatorname{sgn}f_W\bigr).
  \]
  Explain the mapping between signs and geometric names. Produce an illustrated
  table of the eight classes. Test identification after shuffling the solver's
  output order, and explain why the sign of $\det J$ alone is insufficient.

  \item \textbf{Test classification near singularities.}
  Construct short approaches to shoulder, elbow, and wrist singularities.
  Monitor the corresponding factors and the smallest singular value of the
  Jacobian, documenting the scaling used to combine linear and angular quantities.
  Perturb the target poses slightly and compare geometric classification with
  nearest-configuration matching, with and without trajectory history. Examine
  which distinctions disappear or become numerically uncertain. Define tolerances
  for reporting an ambiguous label instead of forcing a sign. At a wrist
  singularity, explain the coupled freedom of the two outer wrist angles.
  Classification can preserve useful identity near a singularity; it cannot
  restore uniqueness when distinct regular branches cease to be distinguishable.

  \item \textbf{Continue each admissible initial class along a path.}
  Prepare three short paths: a regular case, a joint-limit case, and a
  near-singularity case. Specify both position and orientation. From each valid
  starting IK, continue its solution and retain its class history. Compare a
  nearest-solution baseline with selection informed by classification and
  continuation. Do not permit arbitrary jumps between classes: for this PUMA
  model, a continuous change of a separator sign requires a singular crossing.

  \item \textbf{Validate and interpret the trajectories.}
  Compare completed path fraction, joint travel, minimum joint-limit margin,
  singularity proximity, and label consistency. Resolve angle representatives
  consistently with the limits. Refine the path sampling near problematic
  transitions and re-solve the prescribed intermediate poses; feasible endpoints
  alone do not validate a connection. Explain failures and distinguish numerical
  ambiguity, joint-limit violations, and actual singularities. A sampled study
  is evidence, not a proof of continuous-path feasibility.
\end{itemize}

\section*{Available assets and lecture simulations}
\textbf{The PUMA 560 URDF description and STL meshes are provided.}
The URDF is at
\path{lectures_main/assets/models/puma/puma560_robot.urdf};
the corresponding STL files are in the same directory. Paths are relative to
the repository root.

Use the simulations on the lecture website to inspect frames and axes, verify
D--H parameters, visualize IK postures, and simulate or explore paths. Lectures
2 and 4 support modelling and IK; Lectures 5 and 6 support singularities and
classification; Lecture 7 supports path continuation. These materials are in
\path{lectures_main/lectures/}; exercise briefs are in
\path{lectures_main/exercises/}.

The existing Lecture 7 PUMA path demonstration constrains wrist-center position
and holds the wrist joints at their seed values. Use it as a starting point,
then extend or complement it with full-pose IK to study all eight classes along
a position-and-orientation path. Visual inspection should accompany quantitative
checks.

\section*{Expected outcome}
Students should be able to read and interpret the robot's URDF, establish and
validate its D--H parameters, explain wrist-partitioned IK, and identify the eight
regular shoulder--elbow--wrist classes. They should understand how singularities
affect those classes, when identification becomes ambiguous, and why an initial
IK choice matters for motion planning.

Submit reproducible code and test poses/paths, an illustrated classification
table, comparison plots, a concise 5--6-page report, and a short simulation
demonstration. Explain both successful continuations and failures. A positive
result showing that classification always outperforms the baseline is not required.

\section*{Workload and collaboration}
For a standalone 2-ECTS project, budget approximately 60 hours per student
(120 person-hours per pair), following EPFL's course-design convention of
30 hours per credit.\footnote{\href{https://www.epfl.ch/education/teaching/teaching-guide-2/getting-started/design-a-course_1/}{EPFL: Design a Course, workload guidance}.}
If the two credits also include lectures and exercises, deduct their workload
from this allocation. Reuse existing IK and visualization code to keep the
scope manageable; hardware execution and collision-checker development are
outside the core assignment. One student can lead modelling and classification,
and the other continuation and experiments; both must contribute to validation
and be able to explain the complete project.

\end{document}
